\scanpage{42}
% L1
\uline{11}\quad \uline{Interprétation topologique du \deleted{\unsure{d}}\unsure{type|types}} \deleted{\unsure{filtré|filtr.}} \uline{dimensionnel} \unsure{(ou|(Cas} ``géométrique'').

% L3
Soit $X$ \deleted{(projective)} non singulière sur $k$ alg. clos. Considérons la filtration $X$ par la codimension, et la suite spectrale
\[
H^{\unsure{n|*}}(X,\mathbb Q_\ell)\Longleftarrow\inserted{E_1^{pq}=}\coprod_{x\in X[p]} H^{q-p}(x,\mathbb Q_\ell)(-p)\deleted{\unsure{)}}
\]
% slanted
\margin{Attention, terme initial $E_1$ et non $E_2$ !!}
\figure{\begin{tikzpicture}[x=0.7cm,y=0.7cm]
\draw (-0.65,0)--(3.9,0); \draw (0,-0.55)--(0,2.85);
\draw (0,0)--(4,2.9);
\node at (-0.32,2.95) {$q$}; \node at (4.8,-0.18) {$p$};
\node at (4.05,1.08) {$0$};
\end{tikzpicture}}
où on pose
\[
\left\{\begin{aligned}
X[p]&=\{x\in X\mid\dim\mathcal O_{X,x}=p\}\\
H^*(x,\mathbb Z_\ell(-p))&=\underset{\substack{U\text{ ouvert}\\\ne\varnothing\text{ de }\bar x}}{\underrightarrow{\lim}} H^*(U,\mathbb Q_\ell(\unsure{-p|p})))
\end{aligned}\right.
\]

% L11
On veut que \deleted{\unsure{la succession de}} cette suite spectrale provienne d'une suite spectrale de motifs [du moins à partir de $E_r^{pq}$ avec $r\geq2$ ; sinon il faudrait parler en ind-motifs, ou bien prendre une filtration \uline{finie} convenable de $X$ par des sous-schémas fermés]. Le morceau \inserted{en degré $h$} de la filtration $\geq p$ est visiblement de type dimensionnel $\leq h-2p$. On veut que \inserted{si $X$ est \uuline{projectif} (hypothèse essentielle \unsure{vu si|même si} $X$ de dim $1$ !)} \unsure{ce soient|ce sont|ce soit} \uline{exactement} les morceaux de type dimensionnel $h-p$.
% archive sheet number: 41
